\subsection{3 ノード分散(Pi 3 + Pi 4 + Pi 5、12 コア協調)}
Pi 3 も 4 コア SMP 化したので、第 3 のノードは\textbf{4 コア}で分散に参加できる。
これで 3 台合計\textbf{12 コア}(A53$\times$4 + A72$\times$4 + A76$\times$4)の
ヘテロジニアス協調になる。3 台はいずれも本メッシュに参加済みで、計算要求は 2 ノードの
ときと同様に制御面から同時発行し、壁時計を Mac 側で測る。

\medskip\noindent
\textbf{粒度の壁。}ここで新たな制約が効いてくる。各ノードは受け取った列範囲を
\emph{ノード内 4 コアにブロック分割}するため、\textbf{列が 4 本未満のノードでは
1 コアに仕事が偏る}(例：2 列を 4 コアに配ると 1 コアが両列を担い、残り 3 コアは遊ぶ)。
したがって「各ノードに 4 列以上を、かつ重い中央列を最速の Pi 5 に」置くのが要諦になる。
これは第 3 章の\textbf{ノード内}の粒度問題が、\textbf{ノード間}にも二重に現れたものである。
列割当てを振って壁時計を最小化した結果、最良は次の配置であった(表\ref{tab:dist3best})。

\begin{table}[h]\centering
\caption{最良割当ての内訳 ―― $N{=}14$、12 コア、総和 365\,596 \checkmark}\label{tab:dist3best}
\small
\begin{tabular}{lllrr}
\toprule
ノード & コア & 担当列 & 部分解 & 経過 \\
\midrule
Pi 3 (A53) & 4 & [10,14) & 78\,898  &  672\,ms \\
Pi 4 (A72) & 4 & [0,4)   & 78\,898  & 1\,288\,ms \\
Pi 5 (A76) & 4 & [4,10)  & 207\,800 & 1\,456\,ms \\
\midrule
\multicolumn{3}{l}{\textbf{合計 / 壁時計}} & \textbf{365\,596} & \textbf{1\,458\,ms} \\
\bottomrule
\end{tabular}
\end{table}

\noindent
重い中央 6 列 [4,10) を最速の Pi 5 に、軽い端の 4 列ずつを Pi 3・Pi 4 に割り当てた。
Pi 3 は最も非力ながら軽い 4 列を 4 コアで 1 列ずつ捌き\textbf{最速の 672\,ms で完了}、
律速は中央列を持つ Pi 5(1\,456\,ms)である。表\ref{tab:dist-all}に全構成をまとめる。

\begin{table}[h]\centering
\caption{$N{=}14$(解 365\,596)の全構成比較}\label{tab:dist-all}
\begin{tabular}{lrr}
\toprule
構成 & 壁時計 & 対 Pi 5 単独 \\
\midrule
Pi 4 単独(4 コア)                          & 5\,216\,ms & 0.52$\times$ \\
Pi 5 単独(4 コア)                          & 2\,722\,ms & 1.00$\times$(基準) \\
2 ノード(Pi 4+Pi 5、8 コア)                 & 2\,088\,ms & 1.30$\times$ \\
3 ノード(Pi 3 が 1 コア時、9 コア)           & 1\,585\,ms & 1.72$\times$ \\
\textbf{3 ノード(Pi 3 も 4 コア、12 コア)}   & \textbf{1\,458\,ms} & \textbf{1.87$\times$} \\
\bottomrule
\end{tabular}
\end{table}

\noindent
\textbf{読み取り。}12 コア分散は\textbf{1\,458\,ms}で、Pi 5 単独の\textbf{1.87 倍}、
Pi 4 単独の\textbf{3.58 倍}。解の総和も真値に一致する。一方、9 コア(1\,585\,ms)から
12 コアへの改善は\textbf{1.09 倍}にとどまる。理由は明快で、$N{=}14$ の\textbf{列は 14 本
しかなく}、これを 12 コアへ均等に配れない――Pi 3 は 672\,ms で早々に空くが、そこへ
次の(より重い)列を渡すと今度は Pi 3 が律速に回るため、これ以上詰められない。
\textbf{12 コア規模では、律速が「コアの速さ」から「タスクの粒度」へ移る}のである。

\begin{figure}[h]\centering
\begin{tikzpicture}
\begin{axis}[
  width=0.94\linewidth, height=5.2cm,
  xbar, bar width=11pt, xmin=0, xmax=5600,
  xlabel={壁時計(ms)、$N{=}14$、解 365\,596}, xlabel style={font=\small},
  symbolic y coords={{12コア(3×4)},{9コア(4+4+1)},{8コア(4+4)},{Pi5 単独(4)},{Pi4 単独(4)}},
  ytick=data, tick label style={font=\small},
  nodes near coords, nodes near coords style={font=\tiny},
  every node near coord/.append style={anchor=west},
  enlarge y limits=0.16,
]
\addplot coordinates {(5216,{Pi4 単独(4)}) (2722,{Pi5 単独(4)}) (2088,{8コア(4+4)}) (1585,{9コア(4+4+1)}) (1458,{12コア(3×4)})};
\draw[dashed,red!70!black] (axis cs:1172,{Pi4 単独(4)}) -- (axis cs:1172,{12コア(3×4)})
  node[above,black,font=\tiny,rotate=90,anchor=south east]{ヘテロ理想 1172};
\end{axis}
\end{tikzpicture}
\caption{メッシュ分散 N‑Queens($N{=}14$)の構成別・壁時計。台数を増やすほど短縮し、
12 コアで 1\,458\,ms(Pi 5 単独の 1.87$\times$)。赤破線は 3 台の処理能力から定まる
ヘテロ理想(1\,172\,ms)で、実測はその約 80\,\%。}\label{fig:scale}
\end{figure}

\subsection{なぜ 1.87 倍で頭打ちなのか ―― 容量比と粒度ロス}
基準の「Pi 5 単独」は\textbf{すでに 4 コア並列}(A76$\times$4、2\,722\,ms)である。
したがって 1.87 倍とは\textbf{「12 コアが最速ボードの 4 コアの 1.87 倍」}という意味で、
1 コア比ではない。上限が 3 倍にならないのは、12 コアが\textbf{等質でない}からである。
各ボードを 4 コアで全列 $[0,14)$ を解かせて素の処理能力を測ると(表\ref{tab:cap}):

\begin{table}[h]\centering
\caption{各ボードの 4 コア・スループット($N{=}14$、解 365\,596)}\label{tab:cap}
\begin{tabular}{lrr}
\toprule
ボード & 4 コア時間 & 処理速度(解/ms) \\
\midrule
Pi 5 (A76) & 2\,722\,ms & 134.3 \\
Pi 3 (A53) & 3\,402\,ms & 107.5 \\
Pi 4 (A72) & 5\,216\,ms &  70.1 \\
\midrule
\textbf{3 台合計} & --- & \textbf{311.9} \\
\bottomrule
\end{tabular}
\end{table}

\noindent
3 台合計の能力は Pi 5 単独の $311.9/134.3 = \mathbf{2.32}$\textbf{ 倍}にすぎない
(Pi 4・Pi 3 の各コアは A76 より遅く、"Pi 5 換算で約 2.3 台分"である)。ゆえに
\textbf{完全に均等分散できても上限は 2.32 倍(理想時間 $365596/311.9 = 1172$\,ms)}。
実測 1\,458\,ms はこの理想の\textbf{約 80\,\%}($1172/1458$)で、残り 20\,\% が
前述の粒度ロス――最速の Pi 3 が 672\,ms で空くのに重い中央列を持つ Pi 5 が
1\,456\,ms まで走る取りこぼし――である。整理すると:
\[
1.00\times\text{(Pi 5 の 4 コア)}\;\longrightarrow\;
2.32\times\text{(ヘテロ容量の上限)}\;\longrightarrow\;
1.87\times\text{(粒度ロスで実測)}.
\]
律速はハードウェアではなく\textbf{列という粗い分割単位}であり、先頭 2 クイーンでの分割
(約 $14\times14$ タスク)や動的ワークスチールで 2.3 倍近くまで回収できる。

\medskip\noindent
対照的に、\textbf{細粒度で分割できる dining は 3 ボードとも 4.0 倍}(表\ref{tab:smp})を
出しており、仮に分散 dining を組めば 3 ノードで\textbf{約 12 倍}に届く。すなわち
本メッシュの 12 コアは、\textbf{仕事さえ細かく割れれば理想に近い並列度}を持ち、
N‑Queens で頭打ちになるのは\textbf{ハードウェアではなく列という粗い分割単位}に因る。
列をさらに小さな部分木へ割る動的ワークスチールを入れれば、12 コアの実効を
理想へ寄せられる。
